% !TEX root="../../main.tex"

\renewcommand{\thesection}{1.9}

\section{多元多项式}\label{sec:1.9}

\begin{ex}{1}\label{ex:1.9.1}
    设 \(f=x_1+x_3^5-x_1x_2x_3+2x_1^2x_2-x_2^2x_3^2+2x_1x_2+x_1^2x_3+3x_2-4\).
    把 \(f\) 写成齐次分量之和并使每个齐次分量的项呈字典序排列, 求 \(f\) 的次数和首项.
\end{ex}

\begin{solution}
    设 \(f_i\) 表示 \(f\) 的 \(i\) 次齐次分量, 则
    \[
    f=f_0+f_1+f_2+f_3+f_4+f_5,
    \]
    其中
    \begin{gather*}
        f_0=-4,\\
        f_1=x_1+3x_2,\\
        f_2=2x_1x_2,\\
        f_3=2x_1^2x_2+x_1^2x_3-x_1x_2x_3,\\
        f_4=-x_2^2x_3^2,\\
        f_5=x_3^5.
    \end{gather*}

    故 \(f\) 的次数为 \(5\).

    将 \(f\) 按字典序应写为
    \[
    f=2x_1^2x_2+x_1^2x_3-x_1x_2x_3+2x_1x_2+x_1-x_2^2x_3^2+3x_2+x_3^5-4,
    \]
    故 \(\lt(f)=2x_1^2x_2\).
\end{solution}

\begin{ex}{8}\label{ex:1.9.8}
    设 \(f\in\Omega[x_1,x_2,\cdots,x_n]\), \(M\) 表示 \(f\) 中所有项中出现的最大指数,
    且对于任意一组正整数 \(a_1,a_2,\cdots,a_n\leq M+1\), 均有 \(f(a_1,a_2,\cdots,a_n)=0\).
    证明 \(f=0\).
\end{ex}

\begin{proof}
    当 \(n=1\) 时, \(f\) 为一元 \(m\) 次多项式, 其中 \(m\leq M\). 故 \(f\) 至多有 \(m\)个根.
    由于 \(x_1\) 有 \(M+1\) 个取值使得 \(f=0\), 可知 \(f=0\).

    当 \(n=k\) (\(k\in\N^*\)) 时, 假设对满足题意的 \(f\), 有 \(f=0\).
    
    则当 \(n=k+1\), 对于每一组正整数 \(a_1,a_2,\cdots,a_{n-1}\leq M+1\), 设
    \[
    g(x_n)=f(a_1,a_2,\cdots,a_{n-1},x_n)=\sum_{i=1}^{M}f_i(a_1,a_2,\cdots,a_{n-1})x_n^i.
    \]
    由于 \(x_n\) 有 \(M+1\) 个取值使得 \(g(x_n)=0\), 可知 \(g(x_n)=0\). 因此
    \[
    f_i(a_1,a_2,\cdots,a_{n-1})=0,\quad i=1,2,\cdots,M.
    \]
    由归纳假设可知 \(f_i(x_1,x_2,\cdots,x_{n-1})=0\), 故
    \[
    f=\sum_{i=1}^{M}f_i(x_1,x_2,\cdots,x_{n-1})=0.
    \]

    由数学归纳法可知 \(f=0\).
\end{proof}